\documentclass[8pt, xcolor=table,aspectratio=169]{beamer}


\usepackage[utf8]{inputenc}
\usepackage[T1]{fontenc}

\usepackage{helvet}
\usepackage{comment}
\usepackage{bm}

%\usefonttheme[onlymath]{serif}

\usepackage{setspace}
\usepackage{tikz}
\usepackage{subcaption}
\usetikzlibrary{math}
\usetikzlibrary{calc}
\usetikzlibrary{arrows.meta}
\usetikzlibrary{positioning, fit, shapes.geometric, calc}
\usetikzlibrary{decorations.pathreplacing}

\usepackage{adjustbox}

\usepackage{pgfplots}
\pgfplotsset{compat=newest}
\pgfplotsset{every axis/.append style={
      label style={font=\small},
      tick label style={font=\small},
      legend style={font=\small\sf, row sep=-2pt},
      legend cell align={left}
    }
}

\def\centerarc[#1](#2)(#3:#4:#5)% Syntax: [draw options] (center) (initial angle:final angle:radius)
{ \draw[#1] ($(#2)+({#5*cos(#3)},{#5*sin(#3)})$) arc (#3:#4:#5); }

%\usepgfplotslibrary{external}
%\tikzexternalize[prefix=./fig/]
%\tikzexternalenable

\usepackage{amsmath,amssymb,amsfonts,amsthm}
\usepackage{mathtools}

\usepackage[absolute,showboxes,overlay]{textpos} %showboxes,overlay
\TPshowboxestrue            % affiche le contour des textblocks
\TPshowboxesfalse           % fait disparaitre le contour des textblocks
\textblockorigin{2mm}{8mm}  % origine des positions pour placer les textblocks

\usepackage{xcolor,colortbl,color}
\usepackage{framed}
\usepackage{mdframed}

\newtheorem{proposition}{Proposition}[theorem]

\mathtoolsset{showonlyrefs}

\usepackage{graphics}
\usepackage[pdf]{pstricks}
\usepackage{graphicx}
\usepackage{pstricks-add}
\usepackage{pst-node}
\usepackage{pst-plot}
\usepackage{pst-math}
\usepackage{pst-coil}
\usepackage{pst-grad}

% For checking draft
\usepackage{ifdraft}
% Backup slide numbers
\usepackage{appendixnumberbeamer}

% Box
\usepackage{tcolorbox}
\tcbuselibrary{skins, breakable}

%Algorithm
%\usepackage[lined,boxed]{algorithm2e}
\usepackage[boxed]{algorithm2e}
%Multirow
\usepackage{multirow}

%Itemize space
\usepackage{enumitem}
%\setlist[itemize]{label=\textbullet}
\setlist[itemize,1]{label=\scalebox{1.5}{\textbullet}}  % 1.5 = 150% de taille
%\setlist[itemize,2]{label=\scalebox{1.2}{\textbullet}}  % 1.5 = 150% de taille
\setlist[itemize,2]{label=\scalebox{1.2}{\(\circ\)}}
\usepackage{listings}

%\lstset{
%  basicstyle=\ttfamily\scriptsize,
%  keywordstyle=\color{blue}\bfseries,
%  commentstyle=\color{gray}\itshape,
%  language=C++,
%  breaklines=true,
%  frame=single,
%  rulecolor=\color{gray},
%  showstringspaces=false
%}
\lstset{
  basicstyle=\ttfamily\scriptsize,
  keywordstyle=\color{blue}\bfseries,
 escapeinside={(*@}{@*)},
    commentstyle=\color{black}\itshape,
 language=C++,
  breaklines=true,
  frame=single,
  rulecolor=\color{gray},
  showstringspaces=false
}
\setlist[itemize,1]{
    label=$\square$,
    itemsep=0.6em
}

\setlist[itemize,2]{
    label=$\circ$,
    itemsep=0.4em
}

\usepackage{fvextra}

%----------------------------------------------------
% Colors
%----------------------------------------------------

\definecolor{colorField}{rgb}{0.2,0.6,0.2}
\definecolor{colorNumTra}{rgb}{1,0.5,0}
% \definecolor{colorOp}{rgb}{1,0.1,0.1}

\definecolor{darkrose}{rgb}{0.7,0,0.4}
\definecolor{darkblue}{rgb}{0.2,0.2,0.7}
\definecolor{darkorange}{rgb}{0.85,0.3,0.03}

\definecolor{main}{rgb}{0.25,0.35,0.55}
%\definecolor{colorStruct}{rgb}{0.15,0.15,0.45}
%\definecolor{main}{rgb}{0.25,0.35,0.55}
%\definecolor{colorStruct}{rgb}{0.2,0.3,0.5}
\definecolor{colorStruct}{HTML}{19696e}
\definecolor{colorStructLight}{rgb}{0.1,0.2,0.5}

\definecolor{lightlightlightgray}{rgb}{0.93,0.93,0.93}
\definecolor{lightlightgray}{rgb}{0.8,0.8,0.8}
\definecolor{lightgray}{rgb}{0.6,0.6,0.6}
\definecolor{darkgray}{rgb}{0.4,0.4,0.4}
\definecolor{darkdarkgray}{rgb}{0.2,0.2,0.2}

\definecolor{ubuntuorange}{rgb}{1,0.388235294,0.035294118}
\definecolor{ubuntuyellow}{rgb}{1,0.709803922,0.082352941}
\definecolor{ubuntured}{rgb}{0.788235294,0,0.08627451}
\definecolor{lightred}{rgb}{1,0.8,0.8}
\definecolor{lightred}{rgb}{1,0.85,0.7}
\definecolor{darkgreen}{rgb}{0.1,0.5,0.1}
\definecolor{ulggreen}{rgb}{0,0.286274509803922,0.309803921568627}
\definecolor{lightgreen}{rgb}{0.85,0.9,0.85}
\definecolor{lightyellow}{rgb}{1,0.8,0.5}
\definecolor{lightlightyellow}{rgb}{1,0.95,0.7}
\definecolor{bordeaux}{rgb}{0.686274509803922,0.149019607843137,0.149019607843137}
\definecolor{mainTOC}{rgb}{0.788235294,0,0.08627451}
\definecolor{myTeal}{HTML}{008080}
\definecolor{myReddish}{HTML}{d40000}
\definecolor{myGold}{HTML}{d4aa00}

%Used in GPU arch figure
\definecolor{gpumemorycolor}{HTML}{9013FE}
\definecolor{sharedmemorycolor}{HTML}{9B9B9B}

%%%%%%%%%%%%%%%%%%%%%%%%%%%%%%%%%%%%%%%%%%%%%%%%%%%%%%%%%%
\colorlet{colorTrans}{myReddish}
\colorlet{colorVol}{myTeal}
\colorlet{colorOp}{myGold}

\definecolor{lightblue}{rgb}{0.9,0.95,0.98}

%----------------------------------------------------
% Beamer
%----------------------------------------------------

\beamertemplatenavigationsymbolsempty
\setbeamercolor{frametitle}{fg=white,bg=colorStruct}

\newcommand{\axelfootline}{\begin{flushright}{\color{darkgray}\small\insertframenumber/\inserttotalframenumber}\;\;\;\;\end{flushright}}
\setbeamertemplate{footline}{\axelfootline}
\setbeamercolor{button}{bg=main,fg=white}
\setbeamertemplate{itemize items}[bullet]
\setbeamercolor{itemize item}{fg=main}
\setbeamertemplate{itemize subitem}[circle]
\setbeamercolor{itemize subitem}{fg=main}
\setbeamertemplate{itemize/enumerate body begin}{\normalsize}
\setbeamertemplate{itemize/enumerate subbody begin}{\normalsize}
\setbeamerfont{item}{size=\normalsize}
\setbeamerfont{subitem}{size=\normalsize}
\def\myframetitle#1#2{\frametitle{#2 \hspace{0pt plus 1 filll} {\small\color{yellow}[\:#1\:]\!\!\!} }}

\setbeamercolor{block title}{fg=white,bg=colorStruct!70!white}
\setbeamerfont{block title}{size=\normalsize}
\setbeamercolor{block body}{fg=black,bg=lightlightlightgray}
\setbeamertemplate{blocks}[shadow=false]
%\setbeamertemplate{blocks}[rounded][shadow=true]

\setbeamertemplate{itemize items}[circle]
\setbeamertemplate{itemize subitem}{$\circ$}

% ---------
% Symbols
% ---------

% Scalar valued versions
\newcommand{\matalg}[1]{\mathsf{#1}}
\newcommand{\vecalg}[1]{\mathsf{#1}}

% Vector valued versions
\newcommand{\matalgvv}[1]{\mathbf{\mathsf{#1}}}
\newcommand{\vecalgvv}[1]{\mathbf{\mathsf{#1}}}


\newcommand{\prodV}[2]{({\textstyle{#1}},{\textstyle{#2}})}
\newcommand{\prodS}[2]{\langle{\textstyle{#1}},{\textstyle{#2}}\rangle}

\newcommand{\PPs}{\mathbb{P}}
\newcommand{\PPv}{\pmb{\mathbb{P}}}

%% Miscellaneous
\renewcommand{\Re}{\operatorname{Re}}
\renewcommand{\Im}{\operatorname{Im}}

\newcommand{\im}{\imath}
\renewcommand{\i}{\im}

%% Differential operator
\newcommand{\grad}{\boldsymbol \nabla}
\renewcommand{\div}{\grad \cdot}
\newcommand{\curl}{\grad \times}

%% Bold uppercases
\newcommand{\BA}{\mathbf{A}}
\newcommand{\BB}{\mathbf{B}}
\newcommand{\BC}{\mathbf{C}}
\newcommand{\BD}{\mathbf{D}}
\newcommand{\BE}{\mathbf{E}}
\newcommand{\BF}{\mathbf{F}}
\newcommand{\BG}{\mathbf{G}}
\newcommand{\BH}{\mathbf{H}}
\newcommand{\BI}{\mathbf{I}}
\newcommand{\BJ}{\mathbf{J}}
\newcommand{\BK}{\mathbf{K}}
\newcommand{\BL}{\mathbf{L}}
\newcommand{\BM}{\mathbf{M}}
\newcommand{\BN}{\mathbf{N}}
\newcommand{\BO}{\mathbf{O}}
\newcommand{\BP}{\mathbf{P}}
\newcommand{\BQ}{\mathbf{Q}}
\newcommand{\BR}{\mathbf{R}}
\newcommand{\BS}{\mathbf{S}}
\newcommand{\BT}{\mathbf{T}}
\newcommand{\BU}{\mathbf{U}}
\newcommand{\BV}{\mathbf{V}}
\newcommand{\BW}{\mathbf{W}}
\newcommand{\BX}{\mathbf{X}}
\newcommand{\BY}{\mathbf{Y}}
\newcommand{\BZ}{\mathbf{Z}}

%% Bold lowercases
\newcommand{\vectort}[1]{\bm{#1}}
\newcommand{\ba}{\vectort{a}}
\newcommand{\bb}{\vectort{b}}
\newcommand{\bc}{\vectort{c}}
\newcommand{\bd}{\vectort{d}}
\newcommand{\be}{\vectort{e}}
\renewcommand{\bf}{\vectort{f}}
\newcommand{\bg}{\vectort{g}}
\newcommand{\bh}{\vectort{h}}
\newcommand{\bi}{\vectort{i}}
\newcommand{\bj}{\vectort{j}}
\newcommand{\bk}{\vectort{k}}
\newcommand{\bl}{\vectort{l}}
\newcommand{\bbm}{\vectort{m}} % avoid collision with bm
\newcommand{\bn}{\vectort{n}}
\newcommand{\bo}{\vectort{o}}
\newcommand{\bp}{\vectort{p}}
\newcommand{\bq}{\vectort{q}}
\newcommand{\br}{\vectort{r}}
\newcommand{\bs}{\vectort{s}}
\newcommand{\bt}{\vectort{t}}
\newcommand{\bu}{\vectort{u}}
\newcommand{\bv}{\vectort{v}}
\newcommand{\bw}{\vectort{w}}
\newcommand{\bx}{\vectort{x}}
\newcommand{\by}{\vectort{y}}
\newcommand{\bz}{\vectort{z}}
\newcommand{\bzero}{\mathbf{0}}
\newcommand{\bphi}{\boldsymbol{\phi}}

%% Curly uppercases
\newcommand{\CA}{\mathcal{A}}
\newcommand{\CB}{\mathcal{B}}
\newcommand{\CC}{\mathcal{C}}
\newcommand{\CD}{\mathcal{D}}
\newcommand{\CE}{\mathcal{E}}
\newcommand{\CF}{\mathcal{F}}
\newcommand{\CG}{\mathcal{G}}
\newcommand{\CH}{\mathcal{H}}
\newcommand{\CI}{\mathcal{I}}
\newcommand{\CJ}{\mathcal{J}}
\newcommand{\CK}{\mathcal{K}}
\newcommand{\CL}{\mathcal{L}}
\newcommand{\CM}{\mathcal{M}}
\newcommand{\CN}{\mathcal{N}}
\newcommand{\CO}{\mathcal{O}}
\newcommand{\CP}{\mathcal{P}}
\newcommand{\CQ}{\mathcal{Q}}
\newcommand{\CR}{\mathcal{R}}
\newcommand{\CS}{\mathcal{S}}
\newcommand{\CT}{\mathcal{T}}
\newcommand{\CU}{\mathcal{U}}
\newcommand{\CV}{\mathcal{V}}
\newcommand{\CW}{\mathcal{W}}
\newcommand{\CX}{\mathcal{X}}
\newcommand{\CY}{\mathcal{Y}}
\newcommand{\CZ}{\mathcal{Z}}

%% Caligraphic uppercases
\newcommand{\LA}{\mathscr{A}}
\newcommand{\LB}{\mathscr{B}}
\newcommand{\LC}{\mathscr{C}}
\newcommand{\LD}{\mathscr{D}}
\newcommand{\LE}{\mathscr{E}}
\newcommand{\LF}{\mathscr{F}}
\newcommand{\LG}{\mathscr{G}}
\newcommand{\LH}{\mathscr{H}}
\newcommand{\LI}{\mathscr{I}}
\newcommand{\LJ}{\mathscr{J}}
\newcommand{\LK}{\mathscr{K}}
\newcommand{\LL}{\mathscr{L}}
\newcommand{\LM}{\mathscr{M}}
\newcommand{\LN}{\mathscr{N}}
\newcommand{\LO}{\mathscr{O}}
\newcommand{\LP}{\mathscr{P}}
\newcommand{\LQ}{\mathscr{Q}}
\newcommand{\LR}{\mathscr{R}}
\newcommand{\LS}{\mathscr{S}}
\newcommand{\LT}{\mathscr{T}}
\newcommand{\LU}{\mathscr{U}}
\newcommand{\LV}{\mathscr{V}}
\newcommand{\LW}{\mathscr{W}}
\newcommand{\LX}{\mathscr{X}}
\newcommand{\LY}{\mathscr{Y}}
\newcommand{\LZ}{\mathscr{Z}}

%% Caligraphic uppercases
\newcommand{\tp}{\textup{\textrm{p}}}
\newcommand{\TA}{\textup{\textrm{A}}}
\newcommand{\TB}{\textup{\textrm{B}}}
\newcommand{\TC}{\textup{\textrm{C}}}
\newcommand{\TD}{\textup{\textrm{D}}}
\newcommand{\TE}{\textup{\textrm{E}}}
\newcommand{\TF}{\textup{\textrm{F}}}
\newcommand{\TG}{\textup{\textrm{G}}}
%\newcommand{\TH}{\textup{\textrm{H}}}
\newcommand{\TI}{\textup{\textrm{I}}}
\newcommand{\TJ}{\textup{\textrm{J}}}
\newcommand{\TK}{\textup{\textrm{K}}}
\newcommand{\TL}{\textup{\textrm{L}}}
\newcommand{\TM}{\textup{\textrm{M}}}
\newcommand{\TN}{\textup{\textrm{N}}}
\newcommand{\TO}{\textup{\textrm{O}}}
\newcommand{\TP}{\textup{\textrm{P}}}
\newcommand{\TQ}{\textup{\textrm{Q}}}
\newcommand{\TR}{\textup{\textrm{R}}}
\newcommand{\TS}{\textup{\textrm{S}}}
\newcommand{\TT}{\textup{\textrm{T}}}
\newcommand{\TU}{\textup{\textrm{U}}}
\newcommand{\TV}{\textup{\textrm{V}}}
\newcommand{\TW}{\textup{\textrm{W}}}
\newcommand{\TX}{\textup{\textrm{X}}}
\newcommand{\TY}{\textup{\textrm{Y}}}
\newcommand{\TZ}{\textup{\textrm{Z}}}

%% Bold curly uppercases
\newcommand{\BCA}{\boldsymbol{\CA}}
\newcommand{\BCB}{\boldsymbol{\CB}}
\newcommand{\BCC}{\boldsymbol{\CC}}
\newcommand{\BCD}{\boldsymbol{\CD}}
\newcommand{\BCE}{\boldsymbol{\CE}}
\newcommand{\BCF}{\boldsymbol{\CF}}
\newcommand{\BCG}{\boldsymbol{\CG}}
\newcommand{\BCH}{\boldsymbol{\CH}}
\newcommand{\BCI}{\boldsymbol{\CI}}
\newcommand{\BCJ}{\boldsymbol{\CJ}}
\newcommand{\BCK}{\boldsymbol{\CK}}
\newcommand{\BCL}{\boldsymbol{\CL}}
\newcommand{\BCM}{\boldsymbol{\CM}}
\newcommand{\BCN}{\boldsymbol{\CN}}
\newcommand{\BCO}{\boldsymbol{\CO}}
\newcommand{\BCP}{\boldsymbol{\CP}}
\newcommand{\BCQ}{\boldsymbol{\CQ}}
\newcommand{\BCR}{\boldsymbol{\CR}}
\newcommand{\BCS}{\boldsymbol{\CS}}
\newcommand{\BCT}{\boldsymbol{\CT}}
\newcommand{\BCU}{\boldsymbol{\CU}}
\newcommand{\BCV}{\boldsymbol{\CV}}
\newcommand{\BCW}{\boldsymbol{\CW}}
\newcommand{\BCX}{\boldsymbol{\CX}}
\newcommand{\BCY}{\boldsymbol{\CY}}
\newcommand{\BCZ}{\boldsymbol{\CZ}}

%% Bold caligraphic uppercases
\newcommand{\BLA}{\pmb{\LA}}
\newcommand{\BLB}{\pmb{\LB}}
\newcommand{\BLC}{\pmb{\LC}}
\newcommand{\BLD}{\pmb{\LD}}
\newcommand{\BLE}{\pmb{\LE}}
\newcommand{\BLF}{\pmb{\LF}}
\newcommand{\BLG}{\pmb{\LG}}
\newcommand{\BLH}{\pmb{\LH}}
\newcommand{\BLI}{\pmb{\LI}}
\newcommand{\BLJ}{\pmb{\LJ}}
\newcommand{\BLK}{\pmb{\LK}}
\newcommand{\BLL}{\pmb{\LL}}
\newcommand{\BLM}{\pmb{\LM}}
\newcommand{\BLN}{\pmb{\LN}}
\newcommand{\BLO}{\pmb{\LO}}
\newcommand{\BLP}{\pmb{\LP}}
\newcommand{\BLQ}{\pmb{\LQ}}
\newcommand{\BLR}{\pmb{\LR}}
\newcommand{\BLS}{\pmb{\LS}}
\newcommand{\BLT}{\pmb{\LT}}
\newcommand{\BLU}{\pmb{\LU}}
\newcommand{\BLV}{\pmb{\LV}}
\newcommand{\BLW}{\pmb{\LW}}
\newcommand{\BLX}{\pmb{\LX}}
\newcommand{\BLY}{\pmb{\LY}}
\newcommand{\BLZ}{\pmb{\LZ}}

%% Bold text-up
\newcommand{\BTA}{\pmb{\TA}}
\newcommand{\BTB}{\pmb{\TB}}
\newcommand{\BTC}{\pmb{\TC}}
\newcommand{\BTD}{\pmb{\TD}}
\newcommand{\BTE}{\pmb{\TE}}
\newcommand{\BTF}{\pmb{\TF}}
\newcommand{\BTG}{\pmb{\TG}}
%\newcommand{\BTH}{\pmb{\TH}}
\newcommand{\BTI}{\pmb{\TI}}
\newcommand{\BTJ}{\pmb{\TJ}}
\newcommand{\BTK}{\pmb{\TK}}
\newcommand{\BTL}{\pmb{\TL}}
\newcommand{\BTM}{\pmb{\TM}}
\newcommand{\BTN}{\pmb{\TN}}
\newcommand{\BTO}{\pmb{\TO}}
\newcommand{\BTP}{\pmb{\TP}}
\newcommand{\BTQ}{\pmb{\TQ}}
\newcommand{\BTR}{\pmb{\TR}}
\newcommand{\BTS}{\pmb{\TS}}
\newcommand{\BTT}{\pmb{\TT}}
\newcommand{\BTU}{\pmb{\TU}}
\newcommand{\BTV}{\pmb{\TV}}
\newcommand{\BTW}{\pmb{\TW}}
\newcommand{\BTX}{\pmb{\TX}}
\newcommand{\BTY}{\pmb{\TY}}
\newcommand{\BTZ}{\pmb{\TZ}}

%% Matlab style
\newcommand{\MA}{\matalg{A}}
\newcommand{\MB}{\matalg{B}}
\newcommand{\MC}{\matalg{C}}
\newcommand{\MD}{\matalg{D}}
\newcommand{\ME}{\matalg{E}}
\newcommand{\MF}{\matalg{F}}
\newcommand{\MG}{\matalg{G}}
\newcommand{\MH}{\matalg{H}}
\newcommand{\MI}{\matalg{I}}
\newcommand{\MJ}{\matalg{J}}
\newcommand{\MK}{\matalg{K}}
\newcommand{\ML}{\matalg{L}}
\newcommand{\MM}{\matalg{M}}
\newcommand{\MN}{\matalg{N}}
\newcommand{\MO}{\matalg{O}}
\newcommand{\MP}{\matalg{P}}
\newcommand{\MQ}{\matalg{Q}}
\newcommand{\MR}{\matalg{R}}
\newcommand{\MS}{\matalg{S}}
\newcommand{\MT}{\matalg{T}}
\newcommand{\MU}{\matalg{U}}
\newcommand{\MV}{\matalg{V}}
\newcommand{\MW}{\matalg{W}}
\newcommand{\MX}{\matalg{X}}
\newcommand{\MY}{\matalg{Y}}
\newcommand{\MZ}{\matalg{Z}}
\newcommand{\MPi}{\matalg{\Pi}}

% \newcommand{\ma}{\matalg{a}}
\newcommand{\mb}{\matalg{b}}
\newcommand{\mc}{\matalg{c}}
\newcommand{\md}{\matalg{d}}
\newcommand{\me}{\matalg{e}}
\newcommand{\mf}{\matalg{f}}
\newcommand{\mg}{\matalg{g}}
\newcommand{\mh}{\matalg{h}}
% \newcommand{\mi}{\matalg{i}}
% \newcommand{\mj}{\matalg{j}}
% \newcommand{\mk}{\matalg{k}}
% \newcommand{\ml}{\matalg{l}}
% \newcommand{\mm}{\matalg{m}}
% \newcommand{\mn}{\matalg{n}}
% \newcommand{\mo}{\matalg{o}}
\newcommand{\mmp}{\matalg{p}}
\newcommand{\mq}{\matalg{q}}
\newcommand{\mr}{\matalg{r}}
\newcommand{\ms}{\matalg{s}}
\newcommand{\mt}{\matalg{t}}
%\newcommand{\mu}{\matalg{u}}
\newcommand{\mv}{\matalg{v}}
\newcommand{\mw}{\matalg{w}}
\newcommand{\mx}{\matalg{x}}
\newcommand{\my}{\matalg{y}}
\newcommand{\mz}{\matalg{z}}

% Greek
\newcommand{\bpi}{\boldsymbol{\pi}}
\newcommand{\bmu}{\boldsymbol{\mu}}
\newcommand{\bvepsilon}{\boldsymbol{\varepsilon}}
\newcommand{\bxi}{\boldsymbol{\xi}}
\newcommand{\bgamma}{\mathsf{\gamma}}
\newcommand{\BPi}{\mathbf{\Pi}}
\newcommand{\BTPi}{\mathbf{\mathsf{\Pi}}}

% rmsub
\newcommand{\rmsub}[1]{_{\mathrm{#1}}}
\newcommand{\rmsup}[1]{^{\mathrm{#1}}}

\newcommand{\mycolorbar}[3]{%
  \begin{tikzpicture}[rotate=0]
		\begin{axis}[
			hide axis,
			scale only axis,
			height=0pt,
			width=0pt,
			colormap/jet,
			colorbar horizontal,
			point meta min=#1,
			point meta max=#2,
			colorbar style={
				width=#3\linewidth,
				height=2mm,
				rotate=0,
				title=,
				xtick={#1,#2},
				xtick style={draw=none},
				axis line style={draw=none},
				xticklabel style={xshift=0pt,yshift=0pt},
				title style={xshift=2.5cm,yshift=3.2cm},
				%xmode=auto,
			}
			]
		\end{axis}
  \end{tikzpicture}%
}
\newcommand{\surface}[1]{\textcolor{colorTrans}{#1}}
\newcommand{\volume}[1]{\textcolor{myTeal}{#1}}
\definecolor{pythonorange}{RGB}{255,127,0}
\definecolor{pythongreen}{RGB}{0,153,0}
\definecolor{pythonblue}{RGB}{0,0,180}
\captionsetup[figure]{labelformat=empty}
% Define colors for active and inactive sections
%\definecolor{activecolor}{rgb}{1,1,0} % yellow
\definecolor{activecolor}{HTML}{F4D58D}
\definecolor{inactivecolor}{rgb}{0.6, 0.6, 0.6} % equivalent to black!40
\definecolor{rtblue}{RGB}{31,97,163}
\definecolor{rtgreen}{RGB}{46,139,87}
\definecolor{rtred}{RGB}{200,50,50}
\definecolor{rtorange}{RGB}{230,140,30}
\definecolor{rtgray}{RGB}{110,110,110}
%\usefonttheme{serif}
\newcommand{\bol}[1]{\boldsymbol{#1}}
\newcommand{\dd}{\,\mathrm{d}}
\newcommand{\temperature}{\textcolor{pythonblue}{\bol{T}}}
\newcommand{\intensity}{\textcolor{darkgreen}{\bol{I _\nu}}}
\newcommand{\jintensity}{\textcolor{darkgreen}{\bol{J _\nu}}}

\usetikzlibrary{patterns}


\tikzset{
  box/.style={draw,rounded corners=2pt,fill=#1!15,minimum height=0.8cm,align=center,font=\small},
  box/.default=rtblue,
  arr/.style={-{Stealth[length=2.5mm]},thick}
}
\newif\ifbackup
\backupfalse

\newif\ifDGslide
\DGslidetrue

\setbeamertemplate{section in toc}[sections numbered]
%% Initialize a custom counter for tracking the section numbers

% Customize the table of contents spacing
%\setbeamertemplate{section in toc}[] % Remove numbers and add ball for sections
%\setbeamertemplate{subsection in toc}[default] % Default style for subsections

% Adjust spacing between sections and subsections
\setlength{\leftskip}{0pt} % Set left margin for TOC
\setlength{\rightskip}{0pt} % Set right margin for TOC
\setlength{\parskip}{2mm} % Space between paragraphs in TOC (sections and subsections)



\setbeamertemplate{section in toc shaded}[default][100]

% Use \AtBeginSection to customize the TOC slide
\AtBeginSection[]{
    % Change the background color for the TOC slide
    \setbeamercolor{background canvas}{bg=colorStruct} % Set TOC background color
    \setbeamertemplate{footline}{} % Suppress footline (including frame number)
    \begin{frame}[noframenumbering]
      \frametitle{Outline}
      \setbeamercolor{section in toc shaded}{fg=inactivecolor} % Inactive color for shaded sections
      \setbeamercolor{section in toc}{fg=activecolor} % Active color for current section
      \tableofcontents[currentsection, hideothersubsections] % Show current section with hidden subsections
    \end{frame}
    \setbeamertemplate{footline}{\axelfootline} % Reset the footline
    % Reset colors to default after TOC slide
    \setbeamercolor{background canvas}{bg=} % Reset background color
    \setbeamercolor{section in toc}{fg=black} % Reset TOC section color to default
    \setbeamercolor{subsection in toc}{fg=black} % Reset TOC subsection color to default
}

\newcommand{\fluxfigure}{
       \begin{tikzpicture}[xscale=1,yscale=1]
        %\filldraw[color=black, fill=black!10, thick] (0,0) ellipse (1.3 and 0.6);
        \draw[fill=black!5,thick] (0,0) -- (1,0) -- (1.5,0.7) -- (0.5,0.7) -- cycle;
        \draw[thick] (0.5,0.7) -- (1,0);
        \draw[-latex, color=black, thick] (0.6,0.55) -- (0.95,0.80);
        %\draw (0.70,0.12) node {\color{black!50}\footnotesize$K$};
        %\draw (1.03,0.3) node {\color{black!50}\footnotesize$K'$};
        \draw (0.95,0.9) node {\color{black}\scriptsize$\mathbf{n}_{K,F}$};
        \draw (0.70,0.2) node {\scriptsize$F$};
        \draw (0.35,0.15) node {\color{black!75}\footnotesize$K$};
        \draw (1.16,0.53) node {\color{black!75}\footnotesize$K'$};
        \draw[-latex, color=black, thick] (0.80,-0.45) -- (1.50, 0.05);
        \draw[latex-, color=colorTrans, thick] (0.90,-0.65) -- (1.60,-0.15);
        \draw (1.85, 0.17) node {\color{black}$\scriptsize g_{K,F}^{\oplus}$};
        \draw (2.05,-0.75) node
          {\color{colorTrans}$\scriptsize g_{K,F}^{\ominus}\color{black}\ \ (=g_{K',F}^{\oplus})$};
      \end{tikzpicture}
    }

\begin{document}

%%%%%%%%%%%%%%%%%%%%%%%%%%%%%%%%%%%%%%%%%%%%%%%%%%%%%%%%%%%%%%%%%%%%%%%%%%%%%%%%
{\setbeamertemplate{footline}{}
\frame[noframenumbering]{


  \begin{textblock*}{160mm}(-2mm,30mm)
    \textblockcolor{colorStruct}
    \begin{center}
      \vspace{1mm}
      {\LARGE\color{white}{Solving the radiative transfer equations using ray tracing ++ in  \texttt{FreeFem++}}   
}
      \vspace{1mm}
    \end{center}
  \end{textblock*}

  \begin{textblock*}{160mm}(-2mm,50mm)
    \textblockcolor{}
    \centering
    {\large \underline{\textsc{A. Chabib}}, \textsc{P-H. Tournier} } \\ \vspace{0.5mm}
    \textcolor{black!50}{{Laboratoire Jacques-Louis Lions }} \\
    \textcolor{black!50}{{CNRS, Sorbonne Université UPMC, France}}\\
    \textcolor{black!50}{{ALPINES team, Centre Inria de Paris}}

  \end{textblock*}

  \begin{textblock*}{160mm}(-2mm,60mm)
    \textblockcolor{}
  
  \end{textblock*}
  \vspace{1cm}
  \begin{textblock*}{160mm}(-2mm,63mm)
    \textblockcolor{}
    \begin{center}
      %\small
      {\large  \texttt{FreeFem++} days}\\
      \vspace{0.1 cm}
      From 10 to 11th December 2026

    \end{center}
  \end{textblock*}



  \begin{textblock*}{150mm}(25mm,2mm)
    \textblockcolor{}
    \includegraphics[height=16mm]{figures/logos/ljll.png}
    \hspace{0.2cm}
    \includegraphics[height=16mm]{figures/logos/LOGO_CNRS_BLEU.png}
    \hspace{0.2cm}
    \includegraphics[height=15mm]{figures/logos/LOGO_SorbonneU_0.jpg}
    \hspace{0.1cm}
    \includegraphics[height=15mm]{figures/logos/inr_logo_noir.png}
    \hspace{0.2cm}

  \end{textblock*}
  %\begin{textblock*}{117mm}(4mm,100mm)
  %  \textblockcolor{}
  %  \hfill
  %  %\includegraphics[width=25mm]{img/internoise2024.png}
  %\end{textblock*}


%   \begin{textblock*}{150mm}(35mm,58mm)
%     \textblockcolor{}
    
%     \includegraphics[height=18mm]{figures/logos/LOGO_CNRS_BLEU.png}
%     %\hspace{2mm}    
%     %\includegraphics[height=6mm]{figures/logos/logoENSTA.png}
%     %\hspace{2mm}
%     %\includegraphics[height=10mm]{figures/logos/logoIPParis.png}
%   \end{textblock*}
  % \begin{textblock*}{150mm}(0mm,35mm)
  %   \textblockcolor{}
  %   \includegraphics[height=15mm]{figures/logos/labelANR.png}
  % \end{textblock*}
%   \begin{textblock*}{150mm}(55mm,62mm)
%     \textblockcolor{}
%     %\includegraphics[height=8mm]{figures/logos/logoENSTA.png}
%     \includegraphics[height=20mm]{figures/logos/ljll.png}
  
% \end{textblock*}

}} % Title file
%%%%%%%%%%%%%%%%%%%%%%%%%%%%%%%%%%%%%%%%%%%%%%%%%%%%%%%%%%%%%%%%%%%%%%%%%%%%%%%%
\section{Introduction to radiative transfer}
%%%%%%%%%%%%%%%%%%%%%%%%%%%%%%%%%%%%%%%%%%%%%%%%%%%%%%%%%%%%%%%%%%%%%%%%%%%%%%%%
\begin{frame}{\textcolor{activecolor}{Introduction}}
  \begin{columns}[T]
    \begin{column}{0.58\textwidth}
      \begin{itemize}[label=$\square$]\setlength\itemsep{0.6em}
        \item The Sun heats the Earth (short wavelengths); the Earth re-emits in the \textbf{infrared}, which the air partly absorbs and re-emits: this is the \textbf{greenhouse effect}.
        \item \textbf{Goal}: compute the atmospheric temperature $\textcolor{pythonblue}{\temperature (\bol{x})}$ in \textbf{3D}%, accounting for terrain (Chamonix valley), clouds and CO$_2$.
        \item \textbf{Difficulty}: the radiative intensity $\textcolor{darkgreen}{\intensity (\bol{x},\bol{\omega})}$ depends on position (3), direction (2) and frequency (1). The problem has \textbf{7 dimensions}.
        \item \textbf{Motivation}: assess the impact of aviation, where CO$_2$ accounts for only about 1/3 of the warming (the rest comes from NO$_x$ and contrails).
      \end{itemize}
    \end{column}
    \begin{column}{0.38\textwidth}
      \centering
      \begin{block}{Planck's law (black body)}
        \[ B_\nu(\temperature)=\frac{\nu^3}{e^{\nu/\temperature}-1} \]
        {\small Sun: $T_S\approx 5800\,$K\\ Earth: $T_E\approx 288\,$K}
      \end{block}
    \end{column}
  \end{columns}
\end{frame}
%%%%%%%%%%%%%%%%%%%%%%%%%%%%%%%%%%%%%%%%%%%%%%%%%%%%%%%%%%%%%%%%%%%%%%%%%%%%%%%%
\begin{frame}{\textcolor{activecolor}{Radiative transfer equations I}}

  \begin{block}{Radiative transfer along a ray}
    \[
      \bol{\omega}\cdot\nabla \intensity + \kappa_\nu \intensity = \kappa_\nu B_\nu(\temperature),
      \qquad \jintensity := \frac{1}{4\pi}\int_{\mathbb{S}^2} \intensity(\bol{\omega})\dd\bol{\omega}
    \]
  \end{block}
  \medskip 
  \begin{block}{Energy balance of the air}
    \[
      \rho\,\mathbf{u}\cdot\nabla \temperature - \kappa_T\Delta \temperature
      + c_2\!\int_0^\infty\!\kappa_\nu B_\nu(\temperature)\dd\nu
      = c_2\!\int_0^\infty\!\kappa_\nu \jintensity\dd\nu
    \]
  \end{block}
  \medskip
  \begin{itemize}[label=$\square$]\setlength\itemsep{0.8em}
    \item $\kappa_\nu(\bol{x})$: \textbf{absorption} coefficient (air, CO$_2$, clouds); $\rho$: air density; $\mathbf{u}$: wind (Navier--Stokes).
    \item The coupling: the air \textbf{emits} $\kappa_\nu B_\nu(\temperature)$ and \textbf{absorbs} $\kappa_\nu \jintensity$; the difference heats or cools it.
    %\item The general case adds a scattering albedo $a_\nu$ (clouds) and reflective boundary conditions.
  \end{itemize}
\end{frame}
%%%%%%%%%%%%%%%%%%%%%%%%%%%%%%%%%%%%%%%%%%%%%%%%%%%%%%%%%%%%%%%%%%%%%%%%%%%%%%%%
\begin{frame}{\textcolor{activecolor}{Radiative transfer equations II -- eliminating the directions}}
  
  Along a segment, the equation can be solved explicitly (Beer--Lambert law):
  \[
    I(\bol{x},\bol{\omega}) = I(\bol{x}_\Gamma)\,e^{-\int_{[\bol{x},\bol{x}_\Gamma]}\kappa}
    + \int_{[\bol{x},\bol{x}_\Gamma]} e^{-\int_{[\bol{x},\bol{x}']}\kappa}\,S(\bol{x}')\dd\bol{x}'
  \]
  Averaging over $\mathbb{S}^2$, $J$ depends only on $\temperature$ and on the boundary:
  \begin{align*}
    \jintensity[\temperature](\bol{x}) &= \underbrace{\frac{1}{4\pi}\int_\Omega \kappa_\nu B_\nu\big(\temperature(\bol{x}')\big)\frac{e^{-\int_{[\bol{x},\bol{x}']}\kappa_\nu}}{|\bol{x}'-\bol{x}|^2}\dd\bol{x}'}_{\text{Volume integral}}\\
    &\quad+ \underbrace{\frac{1}{4\pi}\int_\Gamma \intensity\Big(\bol{x},\tfrac{\bol{x}'-\bol{x}}{|\bol{x}'-\bol{x}|}\Big)\frac{[(\bol{x}'-\bol{x})\cdot\bol{n}(\bol{x}')]^-}{|\bol{x}'-\bol{x}|^3}e^{-\int_{[\bol{x},\bol{x}']}\kappa_\nu}\dd\Gamma(\bol{x}')}_{\text{Surface integral}}
  \end{align*}
  \alert{ $\Longrightarrow$ We go from 7 to 4 dimensions: $(\bol{x},\nu)$.}
\end{frame}
\section{Numerical scheme}
%%%%%%%%%%%%%%%%%%%%%%%%%%%%%%%%%%%%%%%%%%%%%%%%%%%%%%%%%%%%%%%%%%%%%%%%%%%%%%%%
\begin{frame}{\textcolor{activecolor}{Numerical scheme \small {( PH. Tournier $\&$ O. Pironneau) } }}
  \begin{block}{Fixed-point iteration (convergence and monotonicity proven in \texttt{F. Golse \& O. Pironneau [JCP-2022]})}
    \[
      \mathbf{u}\cdot\nabla \temperature^{n+1}-\kappa_T\Delta \temperature^{n+1}+c_2\!\int_0^\infty\!\kappa_\nu B_\nu(\temperature^{n+1})\dd\nu
      = c_2\!\int_0^\infty\!\kappa_\nu J[\temperature^{n}]_\nu\dd\nu
    \]
  \end{block}
  \begin{columns}[T]
    \begin{column}{0.5\textwidth}
      \textbf{At each iteration $n$}:
      \begin{itemize}[label=$\square$]\setlength\itemsep{0.6em}
        \item compute $J[\temperature^n]_\nu$: volume integral $+$ surface integral;
      \item solve the nonlinear temperature equation (\textbf{Newton}) to get $\temperature^{n+1}$.
      \end{itemize}
    \end{column}
    \begin{column}{0.5\textwidth}
      \textbf{Discretization}: $P^1$ finite elements on a tetrahedral mesh ($N$ vertices).
      \[ \mathbf{j}_\nu = \mathbf{G}_\kappa\,\mathbf{b}_\nu,\quad \mathbf{G}_\kappa\in\mathbb{R}^{N\times N} \]
      \alert{$\mathbf{G}_\kappa$ is dense: $O(N^2)$ memory.}
    \end{column}
  \end{columns}
  %\medskip
  %\textbf{Frequency integral}: the kernels depend on $\nu$ only through $\kappa_\nu$, so frequencies are grouped by value of $\kappa$ (8 values instead of 683 measurements).
\end{frame}
%%%%%%%%%%%%%%%%%%%%%%%%%%%%%%%%%%%%%%%%%%%%%%%%%%%%%%%%%%%%%%%%%%%%%%%%%%%%%%%%
\begin{frame}{\textcolor{activecolor}{Hierarchical matrices in a nutshell}}
  \begin{columns}[T]
    \begin{column}{0.60\textwidth}
      \begin{itemize}[label=$\square$]\setlength\itemsep{0.9em}
        %\item The kernel $e^{-\int\kappa}/|\bol{x}'-\bol{x}|^2$ decays quickly: two \textbf{distant} groups of vertices interact in a simple way.
        \item The kernel $\frac{e^{-\int_{[\bol{x},\bol{x}']}\kappa_\nu}}{|\bol{x}'-\bol{x}|^2}$ decays quickly: two \textbf{distant} groups of vertices interact in a simple way.
        \item Their block is \textbf{low-rank}: $\mathbf{B}\approx\sum_{j=1}^{r}\sigma_j\mathbf{u}_j\mathbf{v}_j^{T}$, $r\ll N$.
        \item The SVD costs $O(N^3)$: we use \textbf{ACA} (Adaptive Cross Approximation), which needs only $\sim 2rN$ coefficients.
        \item A geometric \textbf{cluster tree} partitions the matrix; a block is compressed if $\max(\operatorname{diam}X,\operatorname{diam}Y)\le\eta\operatorname{dist}(X,Y)$.
        \item \textbf{Gain}: memory and matrix-vector product go from $O(N^2)$ to almost linear.
        \item Implementation: \texttt{FreeFEM} + \texttt{Htool} library \small(by P. Marchand \& P-H. Tournier)  .
      \end{itemize}
    \end{column}
    \begin{column}{0.36\textwidth}
      \centering
      \begin{tikzpicture}[scale=0.5]
        \fill[green!45!black!35] (0,0) rectangle (8,8);
        \foreach \i in {0,1,2,3}{
          \fill[red!55] (2*\i,6-2*\i) rectangle (2*\i+2,8-2*\i);
        }
        \draw (0,0) grid[step=2] (8,8);
        \draw[very thick] (0,0) rectangle (8,8);
        \draw[very thick] (4,0) -- (4,8);
        \draw[very thick] (0,4) -- (8,4);
      \end{tikzpicture}

      \medskip
      {\small \textcolor{red!70!black}{red}: dense blocks (near field)\\
      \textcolor{green!40!black}{green}: low-rank blocks (far field)}
    \end{column}
  \end{columns}
\end{frame}
%%%%%%%%%%%%%%%%%%%%%%%%%%%%%%%%%%%%%%%%%%%%%%%%%%%%%%%%%%%%%%%%%%%%%%%%%%%%%%%%
\begin{frame}{\textcolor{activecolor}{Speeding up with ray tracing}}
  \begin{columns}[T]
    \begin{column}{0.58\textwidth}
      \medskip
      The radiation reaching $\bol{x}$ travels along a segment inside the air $\Omega$. If a mountain blocks $[\bol{x},\bol{x}']$, no light goes through it, and the contribution of $\bol{x}'$ is zero (the ground is accounted for by the surface term). So the kernel carries a visibility factor:
      \[
        G^{ij}_\kappa=\frac{1}{4\pi}\int_\Omega
        \underbrace{V(\bol{x}^i,\bol{x}')}_{\in\{0,1\}}\,\kappa\,
        \frac{e^{-\int_{[\bol{x}^i,\bol{x}']}\kappa}}{|\bol{x}^i-\bol{x}'|^2}\,\hat w^j(\bol{x}')\dd\bol{x}'
      \]
      \vspace{-0.4em}
      \begin{itemize}[label=$\square$]\setlength\itemsep{0.6em}
        \item \textbf{Skip} the line integral $\int_{[\bol{x},\bol{x}']}\kappa$ for every hidden pair: many matrix entries are exactly $0$.
        \item \textbf{Early stop}: while tracing a ray, stop once the optical depth $\int\kappa$ is large, since $e^{-\int\kappa}\approx 0$.
        \item $V$ does not depend on $\nu$: compute it \textbf{once} and reuse it for every value of $\kappa$.
      \end{itemize}
    \end{column}
    \begin{column}{0.38\textwidth}
      \centering
      \begin{tikzpicture}[scale=0.8]
        \fill[gray!55] (0,-0.5) -- (0,0.3) -- (2,0.4) -- (3.2,1.0) -- (4,2.2) -- (5,1.0) -- (7.5,0.3) -- (7.5,-0.5) -- cycle;
        \node at (3.9,0.2) {\scriptsize terrain};
        \node at (1.2,3.0) {\scriptsize air $\Omega$};
        \fill (0.8,0.9) circle (2.5pt) node[below left] {\scriptsize $\bol{x}$};
        \fill (2.2,1.8) circle (2.5pt) node[above] {\scriptsize $\bol{x}'$};
        \fill (7,0.9) circle (2.5pt) node[above right] {\scriptsize $\bol{x}''$};
        \draw[green!50!black, very thick] (0.8,0.9) -- (2.2,1.8);
        \draw[red!70!black, very thick, dashed] (0.8,0.9) -- (7,0.9);
        \node[green!40!black, font=\scriptsize, align=center] at (1.0,2.0) {visible};
        \node[red!70!black, font=\scriptsize] at (6.2,1.45) {hidden};
      \end{tikzpicture}

      \smallskip
      {\footnotesize \textcolor{green!40!black}{green}: computed \quad
      \textcolor{red!70!black}{red}: skipped}

      \medskip
      \begin{alertblock}{Caveat}
        \footnotesize $V$ is discontinuous, so blocks that are partly visible may be harder to compress with ACA (higher ranks). Fully hidden blocks are free.
      \end{alertblock}
    \end{column}
  \end{columns}
\end{frame}

%%%%%%%%%%%%%%%%%%%%%%%%%%%%%%%%%%%%%%%%%%%%%%%%%%%%%%%%%%%%%%%%%%%%%%%%%%%%%%%%
\section{Raytracing}
\begin{frame}{\textcolor{activecolor}{What is a ray?}}
 \begin{columns}[T]
    % Left: definition
    \begin{column}{0.42\textwidth}
      %\small
      \vspace{1.2cm}
      A \textbf{ray} is a half-line defined by two quantities:
      \begin{itemize}[label=$\square$]\setlength\itemsep{0.6em}
        \item an \textbf{origin} $\mathbf{o}$ (where it starts),
        \item a \textbf{direction} $\mathbf{d}$ (where it goes, usually a unit vector).
      \end{itemize}
 
      \medskip
      Every point on the ray is given by the parameter $t$:
      \[
        \textcolor{blue!70!black}{P(t)} = \mathbf{o} + t\,\textcolor{red!70!black}{\mathbf{d}}
      \]
      for $t \in [t_{\min}, t_{\max}]$.  
        
      \medskip
      \visible<2->{ 
      \alert{Core question:} what is the first surface hit?
    }
    \end{column}
 
    \begin{column}{0.58\textwidth}
      \vspace{1cm}
  
      \centering
      \resizebox{\textwidth}{!}{\begin{tikzpicture}[>=Stealth, scale=1.0]
 
  % Coordinates
  \coordinate (O) at (0,0);        % ray origin
  \coordinate (H) at (5,2.5);      % hit point
  \coordinate (E) at (7.4,3.7);    % far end (continuation)
  \coordinate (P) at (2.6,1.3);    % generic point on the ray
 

  % Ray: solid up to the hit, dashed afterwards
  %\draw[very thick, blue!70!black, ->] (O) -- (H);
  \draw[thick, blue!70!black, dashed] (O) -- (H);
  \visible<1>{\draw[thick, blue!70!black, dashed] (H) -- (E);}
 
  % Direction vector d
  \draw[very thick, red!70!black, -{Stealth[length=3mm]}] (O) -- (1.6,0.8)
    node[midway, below right, font=\small] {$\mathbf{d}$};
  % P(t)
  \fill (P) circle (0pt)   node[above left, blue!70!black, font=\small] {$P(t)$};

% Origin 
  \fill (O) circle (2.2pt) node[below left] {$\mathbf{o}$};

\visible<2->{ 
  % Surface (seen from the side)
  \draw[very thick, gray!70!black] ($(H)+(-0.45,-1.5)$) -- ($(H)+(0.45,1.5)$)
    node[above, font=\footnotesize, text=gray!70!black] {Surface};
 
  % Normal at the hit point
  \draw[thick, green!50!black, -{Stealth}] (H) -- ($(H)+(-1.25,0.38)$)
    node[above left, font=\small] {$\mathbf{n}$};


  % Points
  \fill[orange!90!black] (H) circle (2.5pt)
    node[below right, font=\small] {$P(t_{\mathrm{hit}})$};
 
  % Brace showing the parameter t_hit
  \draw[decorate, decoration={brace, amplitude=6pt, mirror}, thick]
    (O) -- (H) node[midway, below=8pt, font=\small] {$t_{\mathrm{hit}}$};
 }
 
\end{tikzpicture}}
    \end{column}
  \end{columns}
\end{frame}
%%%%%%%%%%%%%%%%%%%%%%%%%%%%%%%%%%%%%%%%%%%%%%%%%%%%%%%%%%%%%%%%%%%%%%%%%%%%%%%%

\begin{frame}{\textcolor{activecolor}{The quadratic wall}}
  \medskip
  \begin{itemize}[label=$\square$]\setlength\itemsep{0.6em}
    \item $V$ is symmetric with a zero diagonal: only pairs $i<j$ matter.
    \[ \text{numPairs}=\frac{N(N-1)}{2} \]
    \item Each pair = one \alert{ray}, each ray must be tested against the obstacles.
  \end{itemize}
  \vspace{2mm}
  
  %\begin{alertblock}{Key ideas:}
    \begin{itemize}[label=$\square$]\setlength\itemsep{0.6em}
      \item Store the result in a \alert{triangular bitset}.
      \begin{center}
  \begin{tabular}{|rrrr|}
    \hline
    $N$ (triangles) & pairs & bitset (1 bit/pair) & full \texttt{long} matrix \\ \hline
    $10^3$ & $5\cdot10^5$ & 62 KB & 8 MB \\
    \hline
    $10^4$ & $5\cdot10^7$ & 6 MB & 800 MB \\
        \hline
    $10^5$ & $5\cdot10^9$ & 625 MB & 80 GB \\ \hline
  \end{tabular}
  \end{center}

    %\item Intersection detection is computationally expensive $\Longrightarrow$ Use the GPU's \alert{ray-tracing hardware} to test the rays.
    \item  Testing millions or even billions of ray–mesh intersections is computationally expensive, even on a GPU.
    $\Longrightarrow$ Reducing the number of intersection tests with \textbf{Bounding Volume Hierarchy (BVH)}
  \end{itemize}
  %\end{alertblock}
\end{frame}
%%%%%%%%%%%%%%%%%%%%%%%%%%%%%%%%%%%%%%%%%%%%%%%%%%%%%%%%%%%%%%%%%%%%%%%%%%%%%%%%
\begin{frame}{\textcolor{activecolor}{BVH: a hierarchy of bounding boxes}}
  \begin{columns}[T]
    \begin{column}{0.52\textwidth}
      \vspace{1cm}
        \begin{center}
  \begin{itemize}[label=$\square$]\setlength\itemsep{0.6em}
    \item Each node = an axis-aligned box (AABB) containing its whole subtree.
    \item The \alert{leaves} hold a few triangles.
    \item Building the BVH is expensive, the cost is amortized over millions of rays
    \item If a ray misses a box, it misses \alert{everything inside it}.
      \begin{itemize}[label=$\square$]\setlength\itemsep{0.6em}
      \item Test the ray against the root box (ray-box test, very cheap).
      \item If hit: descend into both children; otherwise \alert{prune} the whole subtree.
      \item On reaching a leaf: test the ray against its point.
      \item Stop at the first hit.
      \end{itemize}
 \end{itemize}
  \end{center}
    \end{column}
    \begin{column}{0.48\textwidth}
      \medskip
            \centering
      \begin{tikzpicture}[scale=0.75]
        % triangles (small dots)
        \foreach \x/\y in {0.5/0.5,1.1/0.9,0.6/1.5,1.4/1.7}{\fill[red!70] (\x,\y) circle (3pt);}
        \foreach \x/\y in {3.3/0.4,3.9/0.9,3.5/1.4,4.3/0.6}{\fill[red!70] (\x,\y) circle (3pt);}
        \foreach \x/\y in {1.8/3.3,2.5/3.7,3.2/3.2,2.6/2.8}{\fill[red!70] (\x,\y) circle (3pt);}
        % leaf boxes
        \draw[green,thick] (0.2,0.2) rectangle (1.7,2.0);
        \draw[green,thick] (3.0,0.1) rectangle (4.6,1.7);
        \draw[green,thick] (1.5,2.5) rectangle (3.5,4.0);
        % intermediate box
        \draw[orange,thick,dashed] (0.0,0.0) rectangle (4.8,2.2);
        % root
        \draw[blue,very thick] (-0.2,-0.2) rectangle (5.0,4.2);
        \node[blue,font=\scriptsize] at (4.3,3.8) {Root};
      \end{tikzpicture}

      \vspace{0.5 cm}
      
      \begin{tikzpicture}[level distance=1.2cm,
          level 1/.style={sibling distance=2.4cm},
          level 2/.style={sibling distance=1.3cm},
          every node/.style={font=\scriptsize}]
        \node[box=blue] {Root}
          child { node[box=orange] {Box A}
            child { node[box=green] {L1} }
            child { node[box=green] {L2} } }
          child { node[box=green] {L3} };
      \end{tikzpicture}
    \end{column}
  \end{columns}


\end{frame}
%%%%%%%%%%%%%%%%%%%%%%%%%%%%%%%%%%%%%%%%%%%%%%%%%%%%%%%%%%%%%%%%%%%%%%%%%%%%%%%%
%\begin{frame}{\textcolor{activecolor}{The hardware side: RT Cores}}
%  \begin{columns}[T]
%    \begin{column}{0.52\textwidth}
%      \vspace{1.8cm}
%      On an NVIDIA RTX GPU, each SM contains, besides the regular CUDA cores, an \alert{RT Core}:
%      \begin{itemize}[label=$\square$]\setlength\itemsep{0.6em}
%        \item A \emph{fixed-function} unit that walks the BVH.
%        \item Performs ray--box and ray--triangle tests in hardware.
%        \item CUDA cores only \alert{launch} the ray and read back the result.
%        \item Compilation target here: \texttt{sm\_86} (Ampere, 2nd-generation RT Cores).
%        \item The repetitive and costly part (tree walking, billions of box tests) is delegated to a dedicated circuit, while the SMs keep computing.
%        
%      \end{itemize}
%    \end{column}
%
%    \begin{column}{0.48\textwidth}
%      \vspace{1.5cm}
%      \centering
%      \begin{tikzpicture}[node distance=0.9cm]
%        \node[box=blue,minimum width=3.6cm] (sm) {Shader on CUDA cores\\(\texttt{raygen})};
%        \node[box=orange,minimum width=3.6cm,below=1.3cm of sm] (rt) {RT Core\\BVH traversal + intersection};
%        \node[box=gray,minimum width=3.6cm,below=1.3cm of rt] (mem) {BVH in GPU memory};
%        \draw[arr] ([xshift=-0.6cm]sm.south) -- node[left,font=\scriptsize] {ray} ([xshift=-0.6cm]rt.north);
%        \draw[arr] ([xshift=0.6cm]rt.north) -- node[right,font=\scriptsize] {hit / miss} ([xshift=0.6cm]sm.south);
%        \draw[arr,<->] (rt) -- (mem);
%      \end{tikzpicture}
%    \end{column}
%  \end{columns}
%  \vspace{2mm}
%
%\end{frame}
%%%%%%%%%%%%%%%%%%%%%%%%%%%%%%%%%%%%%%%%%%%%%%%%%%%%%%%%%%%%%%%%%%%%%%%%%%%%%%%%
\begin{frame}{\textcolor{activecolor}{The hardware side: RT Cores}}
  \begin{columns}[T]
    \begin{column}{0.52\textwidth}
      \vspace{1.8cm}
      Modern NVIDIA GPUs (RTX) contain a \alert{dedicated circuit} for ray tracing, the \alert{RT Core}:
      \begin{itemize}[label=$\square$]\setlength\itemsep{0.8em}
        \item The regular GPU cores are \emph{general-purpose}: they run any program, in parallel.
        \item The RT Core does \alert{one thing only}: given a ray, it finds the first object hit, by descending the tree of nested boxes (the BVH).
        \item It is the same operation repeated billions of times: ideal to wire directly into silicon.
        \item The general-purpose cores just \alert{send} the ray and \alert{read} the answer; meanwhile they stay free for other computations.
      \end{itemize}
    \end{column}

    \begin{column}{0.48\textwidth}
      \vspace{1.5cm}
      \centering
      \begin{tikzpicture}[node distance=0.9cm]
        \node[box=blue,minimum width=3.6cm] (sm) {General-purpose cores\\(our program)};
        \node[box=orange,minimum width=3.6cm,below=1.3cm of sm] (rt) {RT Core\\finds the first object hit};
        \node[box=gray,minimum width=3.6cm,below=1.3cm of rt] (mem) {Scene tree (BVH) in memory};
        \draw[arr] ([xshift=-0.6cm]sm.south) -- node[left,font=\scriptsize] {ray} ([xshift=-0.6cm]rt.north);
        \draw[arr] ([xshift=0.6cm]rt.north) -- node[right,font=\scriptsize] {hit / miss} ([xshift=0.6cm]sm.south);
        \draw[arr,<->] (rt) -- (mem);
      \end{tikzpicture}
    \end{column}
  \end{columns}
\end{frame}
%%%%%%%%%%%%%%%%%%%%%%%%%%%%%%%%%%%%%%%%%%%%%%%%%%%%%%%%%%%%%%%%%%%%%%%%%%%%%%%%
\begin{frame}{\textcolor{activecolor}{OptiX: the API driving the RT Cores}}
  \medskip
  OptiX is a \alert{framework with blanks}: NVIDIA provides the tree and its traversal, we only write
  small programs saying what to do \alert{before} and \alert{after} a ray is traced.
  \vspace{2mm}
  \begin{center}
  \begin{tikzpicture}[node distance=1.2cm and 1.4cm]
    \node[box=green] (rg) {Ray launch\\\scriptsize \texttt{raygen}};
    \node[box=blue,right=of rg] (gas) {Scene tree (BVH)\\\scriptsize traversed by the RT Core};
    \node[box=red,above right=0.3cm and 1.4cm of gas] (hit) {If hit\\\scriptsize \texttt{hit group} (empty here)};
    \node[box=orange,below right=0.3cm and 1.4cm of gas] (ms) {If no hit\\\scriptsize \texttt{miss}};
    \draw[arr] (rg) -- node[above,font=\scriptsize] {ray} (gas);
    \draw[arr] (gas) -- (hit);
    \draw[arr] (gas) -- (ms);
  \end{tikzpicture}
  \end{center}
  \medskip
  \begin{itemize}[label=$\square$]\setlength\itemsep{0.6em}
    \item \textbf{We write}: where each ray starts and what to do on a hit or a miss.
    \item \textbf{NVIDIA provides}: the BVH over our triangles (the \emph{GAS}) and its traversal.
    \item \textbf{Bookkeeping} (pipeline, SBT): assembles the programs and tells OptiX which one to call; trivial with a single object.
  \end{itemize}
\end{frame}
\section{Implementation}
%%%%%%%%%%%%%%%%%%%%%%%%%%%%%%%%%%%%%%%%%%%%%%%%%%%%%%%%%%%%%%%%%%%%%%%%%%%%%%%%
\begin{frame}{\textcolor{activecolor}{Pipeline overview}}
  \begin{center}
  \resizebox{\textwidth}{!}{%
  \begin{tikzpicture}[node distance=0.7cm,font=\small]
    \node[box=blue,text width=3.2cm]   (a) {\textbf{1. Prepare}\\[1pt]\scriptsize points, normals, obstacle triangles from the FreeFEM meshes};
    \node[box=orange,text width=3.2cm,right=of a] (b) {\textbf{2. Set up}\\[1pt]\scriptsize copy to GPU, build the BVH, load the programs (once)};
    \node[box=orange,text width=3.2cm,right=of b] (c) {\textbf{3. Trace}\\[1pt]\scriptsize one ray per pair $(i,j)$, in batches};
    \node[box=blue,text width=3.2cm,right=of c]   (d) {\textbf{4. Collect}\\[1pt]\scriptsize bits $\rightarrow$ matrix $V$};
    \draw[arr] (a)--(b);
    \draw[arr] (b)--(c);
    \draw[arr] (c)--(d);
    \node[below=0.15cm of a,blue]   {CPU};
    \node[below=0.15cm of b,orange] {GPU};
    \node[below=0.15cm of c,orange] {GPU};
    \node[below=0.15cm of d,blue]   {CPU};
  \end{tikzpicture}}
  \end{center}
  \vspace{3mm}
  \begin{itemize}[label=$\square$]\setlength\itemsep{0.6em}
    \item \alert{Input}: points $x_1,\dots,x_N$ and the obstacle mesh.
    \item \alert{Output}: the matrix $V\in\{0,1\}^{N\times N}$, with $V_{ij}=1$ iff the segment $[x_i,x_j]$ meets no obstacle.
    \item \alert{Cost}: $\displaystyle \frac{N(N-1)}{2}$ independent tests; only step 3 scales with $N^2$, and it runs on the RT Cores.
    \item Called from \texttt{FreeFEM++} as \texttt{OptixBuildVisibilityPairs3D(Th3, ThS, V)}.
    \item The \texttt{.cu} is compiled separately to .ptx (\texttt{nvcc -ptx}) and loaded at runtime
  \end{itemize}
\end{frame}
%%%%%%%%%%%%%%%%%%%%%%%%%%%%%%%%%%%%%%%%%%%%%%%%%%%%%%%%%%%%%%%%%%%%%%%%%%%%%%%%
\begin{frame}[fragile]{\textcolor{activecolor}{The raygen program}}
\begin{center}
  
  \begin{lstlisting}
extern "C" __global__ void __raygen__rg() {
  const unsigned long long k = optixGetLaunchIndex().x;
  if (k >= params.numPairs) return;
  unsigned a, b;  unrank(k, params.N, a, b);       // k -> (i,j) global idx

  float3 P = params.pts[a], Q = params.pts[b];
  float3 d = Q - P;  float L = sqrtf(dot3(d,d));   // segment length
  d /= L;                                          // unit direction

  unsigned occluded = 1;                           // pessimistic
  optixTrace(params.gas, P, d,
             params.tmin, params.tmax, 0.0f,   // tmin, tmax, time
             OptixVisibilityMask(255),
             TERMINATE_ON_FIRST_HIT | DISABLE_CLOSESTHIT,
             0, 1, 0, occluded);
  if (!occluded) atomicOr(&params.bits[k >> 5], 1u << (k & 31));
}
\end{lstlisting}
\end{center}
 \begin{itemize}
        \item \alert{Interval} $[t_{\min},t_{\max}]=[10^{-6}L,\ (1-10^{-5})L]$: the ray starts \alert{just after} $P$ and stops \alert{just before} $Q$.
        \item \texttt{TERMINATE\_ON\_FIRST\_HIT}: a single obstacle is enough to conclude ``blocked''.
        \item \texttt{DISABLE\_CLOSESTHIT}: we do not compute the closest hit (useless here).
        \item \alert{Payload} = 1 integer (\texttt{occluded}) shared between raygen and miss.
        \item Mask 255: the ray sees all the geometry.
  \end{itemize}
\end{frame}
%%%%%%%%%%%%%%%%%%%%%%%%%%%%%%%%%%%%%%%%%%%%%%%%%%%%%%%%%%%%%%%%%%%%%%%%%%%%%%%%

\end{document}